\Needspace{13\baselineskip}
\section{Companion lab: interpret a bytecode program}
\label{revision-lab:bytecode}

The expression evaluator reads characters whose grammar determines the next
operation. A bytecode interpreter reads an already encoded sequence of operations.
The bytes still need validation, and their effects belong to a second machine
that your MIPS program maintains. This companion revises original project
\texttt{018-tiny-bytecode-interpreter}.

Import \texttt{018-bytecode-machine}, version 0.1.0, through the MARS Lessons
window. Begin with the complete \texttt{worked/record.asm}, then compare
\texttt{worked/truncated.asm}. Implement \texttt{run\_vm} in
\texttt{independent/vm.asm}; replace its single \texttt{break}, keeping the supplied
caller and data arenas. Private labels, frames and dispatch design are yours to
choose. The original project remains under
\texttt{packages/018-tiny-bytecode-interpreter}.

The preceding nested-data chapter supplies bounds and validation reasoning; the
expression companion supplies checked-word arithmetic. Review functions and stack
preservation as needed. Inline predictions connect the written reasoning to
source pointers. Guided Debug first inspects the complete record classifier,
then the independent VM after you have implemented it. You may also work freely
and return to the staged assessments.

\subsection*{Keep the guest machine distinct from the executing MIPS program}

The MIPS program counter selects instructions such as \texttt{lbu},
\texttt{beqzc} and \texttt{sw}. A \emph{guest PC} is an integer byte offset into
the supplied bytecode image. A guest opcode chooses which native handler to run.
For a guest branch to offset four, your interpreter updates that integer and
returns to its dispatch loop. It does not jump the native machine into the
program's data bytes.

The guest operand stack is another caller-owned data array. The native stack
addressed by \texttt{sp} stores your saved registers and private work. A guest
DUP adds a value to the operand array; it does not duplicate a native return
address. Keep guest depth and native frame size separate even if both happen
to be small integers.

\begin{table}[htbp]
  \centering
  \begin{tabular}{r l l}
    \toprule
    Byte offset & Field                              & Ownership                  \\
    \midrule
    0           & Guest PC                           & Output word                \\
    4           & Live stack depth                   & Output word                \\
    8           & Retired guest count                & Output word                \\
    12          & Status                             & Output word                \\
    16          & Capacity, zero through eight words & Preserved input word       \\
    20          & Fault offset, or $-1$              & Output word                \\
    24          & Full guest-stack address           & Preserved eight-byte input \\
    \bottomrule
  \end{tabular}
  \caption{The eight-byte-aligned descriptor occupies 32 bytes. Output fields report the new invocation; they do not request resumption. The word-aligned guest stack is a separate caller-owned subview. A full stack pointer does not fit in one word.}
  \label{tab:revision-bytecode-descriptor}
\end{table}


The public arguments are full code address, selected length, full descriptor
address and guest instruction budget in \texttt{a0} through \texttt{a3}.
Length and budget each lie between zero and 64. Code is a byte slice inside a
96-byte arena and may be unaligned. The descriptor's input pointer names a
word-aligned stack subview with capacity zero through eight words inside a
sixteen-word allocation. Only the selected capacity is available. Extra
physical cells are not permission to push, read ahead or clear a larger array.

Every call starts at guest PC zero, depth zero and retired count zero. Ignore
the old output fields, even if they describe a previous budget return. The
caller may replace the image at the same address and length. This interface
reports a bounded execution; it does not resume that report on the next call.

\subsection*{Decode instruction boundaries before accepting branches}

\begin{table}[htbp]
\centering
\begin{tabular}{l r l}
\toprule
Opcode and name & Bytes & Guest action\\
\midrule
\texttt{00 HALT} & 1 & Advance, retire and halt\\
\texttt{01 PUSH8} & 2 & Push one signed-byte immediate\\
\texttt{02 PUSH32} & 5 & Push one little-endian signed-word immediate\\
\texttt{10 ADD}, \texttt{11 SUB} & 1 & Replace older left and top right by result\\
\texttt{12 MUL}, \texttt{13 DIV} & 1 & Replace older left and top right by result\\
\texttt{20 DUP}, \texttt{21 DROP} & 1 & Copy or remove top\\
\texttt{22 SWAP}, \texttt{23 OVER} & 1 & Exchange top pair or copy second from top\\
\texttt{30 JMP}, \texttt{31 JZ} & 2 & Absolute byte target; JZ always pops condition\\
\bottomrule
\end{tabular}
\caption{Opcodes are hexadecimal; byte sizes are decimal. Arithmetic uses checked signed words. Branch operands are unsigned byte offsets from the code base. They are neither native instruction addresses nor relative displacements.}
\label{tab:revision-bytecode-encoding}
\end{table}


Walk the entire image from offset zero. At each established instruction start,
read the opcode, determine its size and prove that the complete record fits
before advancing. PUSH8's immediate is a signed byte; PUSH32 carries four
little-endian bytes representing any signed-word bit pattern. Unlike the
expression language's positive decimal literal, PUSH32 can encode minus
2147483648 directly as \texttt{02 00 00 00 80}.

The complete \texttt{worked/record.asm} isolates that first decision. Its
\texttt{record\_width} receives a full pointer, selected length and candidate byte
offset in \texttt{a0/a1/a2}. It proves that the opcode byte exists, classifies its
width, then compares the required width with the remaining selected bytes. On
success it returns status zero and width; on failure it returns a local status
and zero width. This helper does not decode an operand, mark a whole image,
validate a branch target or execute a guest instruction.

Debug the five-byte PUSH32 example. Before \texttt{width\_call}, expect selected
length five and candidate offset zero. At \texttt{record\_fits}, before
\texttt{bltu}, \texttt{t1} holds opcode two, \texttt{t2} holds width five and
\texttt{t3} holds remaining length five. At \texttt{width\_ready}, expect
\texttt{v0}=0 and \texttt{v1}=5. Stop Debug, then open
\texttt{worked/truncated.asm}. Its initializer still contains five physical
bytes, but the selected length is four. Follow the opcode read, then the same
width comparison: required five exceeds remaining four. The returned pair is
\texttt{(-2,0)}. Physical storage past the selected end does not complete the record.

Run the ordinary complete example afterward. It stores status and width and
exits silently. Use \texttt{worked/record.kts} to check all opcode values, each at
selected lengths zero through five, plus the complete caller. These local checks support your
first structural pass; they are not acceptance of the VM.

Record every instruction start. With at most 64 code bytes, a doubleword bitmap
can use bit $p$ for start offset $p$. A private array of 64 byte markers is
another choice. A bitmap needs a doubleword shift, including at offset 63;
a word shift would alias upper-half positions. Do not infer a start merely
because an operand byte happens to equal an opcode.

Only after this structural pass succeeds, scan all encoded branches in image
order. Their unsigned target byte must be below the selected length and belong
to the decoded start set. A target equal to length is invalid. Check unreachable
branches and untaken conditional targets too. A valid runtime path does not
repair an invalid encoded program.

Structural failure takes priority over every bad target. For the bytes
\texttt{30 01 ff}, the first record is a JMP with a bad target into its operand,
but the structural pass next finds invalid opcode \texttt{ff} at offset two.
Report invalid opcode there; target checking never begins. Similarly,
\texttt{00 ff} fails validation after its HALT byte. No guest instruction has
executed, and no guest budget has been consumed in either case.

\Needspace{25\baselineskip}
\begin{table}[!ht]
\centering
\begin{tabular}{r l l}
\toprule
Byte offset & Encoded bytes & Instruction\\
\midrule
0 & \texttt{01 00} & PUSH8 0\\
2 & \texttt{01 03} & PUSH8 3\\
4 & \texttt{20} & DUP\\
5 & \texttt{31 10} & JZ 16\\
7 & \texttt{22} & SWAP\\
8 & \texttt{23} & OVER\\
9 & \texttt{10} & ADD\\
10 & \texttt{22} & SWAP\\
11 & \texttt{01 01} & PUSH8 1\\
13 & \texttt{11} & SUB\\
14 & \texttt{30 04} & JMP 4\\
16 & \texttt{21} & DROP\\
17 & \texttt{00} & HALT\\
\bottomrule
\end{tabular}
\caption{The 18-byte countdown image sums three, two and one. The instruction at offset five has target byte hexadecimal \texttt{10}, decimal 16. Offset six is its operand, even though its byte value is also the ADD opcode. Instruction identity comes from decoding at an established start.}
\label{tab:revision-bytecode-countdown}
\end{table}


Change the target operand at offset six of this image from \texttt{10} to
\texttt{06}. Six is inside the image, but not an instruction start. Even at
zero guest budget, validation must return bad-target status minus three, fault
offset five, PC/depth/retired zero. The fault identifies the branch instruction,
not the operand's location or its target value. This distinction is a useful
check against a range-only validator.

\begin{figure}[!ht]
\centering
\begin{tikzpicture}[alt={Actual native descriptor after zero-budget whole-image validation rejects a branch into an operand. The first three output words are zero, status is minus three and fault is five. Native outlines identify the zero execution prefix and the fault field.}]
\node[anchor=north west,inner sep=0,font=\small\bfseries] (title) at (0,0) {Zero budget still requires a valid target};
\node[anchor=north west,inner sep=0] (a) at ([yshift=-2mm]title.south west) {\MARSRevisionCrop{revision/captures/bytecode/vm-bad-target.png}{0.99\linewidth}{0.03125}{0.31555556}{0.58203125}{0.60888889}};
\begin{scope}[shift={(a.south west)},x={($(a.south east)-(a.south west)$)},y={($(a.north west)-(a.south west)$)}]
\draw[black,line width=.7pt] (.20,.3088) rectangle (.77,.6176);\draw[black,line width=.7pt] (.39,.005) rectangle (.58,.303);
\end{scope}
\end{tikzpicture}
\caption{The image targets operand byte six from the JZ at offset five. PC, depth and retired count remain zero; status is minus three and fault five. The fault names the invalid branch instruction, not its target. No guest instruction has retired. Native LaTeX supplies the title, crop and outlines.}
\label{fig:revision-bytecode-native-bad-target}
\end{figure}


\subsection*{Trace a guest loop through stack effects}

The countdown starts by pushing accumulator zero and counter three. Write the
live stack oldest to newest: \([0,3]\). At offset four, DUP copies the counter,
giving \([0,3,3]\). JZ then consumes the copied condition on either path.
For a nonzero counter, execution falls through to seven with \([0,3]\).
The original counter is still present because DUP made the condition copy.

SWAP gives \([3,0]\), and OVER copies the second value from the top, giving
\([3,0,3]\). ADD consumes the older left operand zero and top right operand
three, producing \([3,3]\). After SWAP, PUSH8 1 and SUB, the stack is
\([3,2]\): accumulated sum three and next counter two. JMP selects offset four.
The next iterations return there with \([5,1]\), then \([6,0]\).

The next DUP/JZ pair takes the exit, leaving \([6,0]\). DROP removes the old
counter, and HALT finishes with the single live value six. A native Memory row
can display inactive words after that live value. Their contents are not more
guest results: the descriptor's depth determines which prefix is meaningful.
The valid countdown needs room for three simultaneous values even though its
final result has depth one.

\begin{figure}[!ht]
\centering
\begin{tikzpicture}[alt={Actual native descriptor after countdown halt. The first row is PC eighteen, depth one, retired thirty-three and status one. The second row retains capacity eight and the full stack pointer and reports fault minus one. Native outline identifies the four main output words.}]
\node[anchor=north west,inner sep=0,font=\small\bfseries] (title) at (0,0) {Halt follows thirty-three guest retirements};
\node[anchor=north west,inner sep=0] (a) at ([yshift=-2mm]title.south west) {\MARSRevisionCrop{revision/captures/bytecode/vm-halted.png}{0.99\linewidth}{0.03125}{0.31555556}{0.58203125}{0.60888889}};
\begin{scope}[shift={(a.south west)},x={($(a.south east)-(a.south west)$)},y={($(a.north west)-(a.south west)$)}]
\draw[black,line width=.7pt] (.20,.3088) rectangle (.965,.6176);
\end{scope}
\end{tikzpicture}
\caption{The first row reports PC 18, depth one, retired count 33 and halt status one. Fault is minus one in the second row. The single live result six is verified separately; the stack cells are not pictured here. Native LaTeX supplies the title, crop and outlines.}
\label{fig:revision-bytecode-native-halted}
\end{figure}


\begin{table}[htbp]
\centering
\begin{tabular}{r r r l l}
\toprule
Guest budget & PC & Retired & Live stack, oldest first & Outcome\\
\midrule
0 & 0 & 0 & Empty & Budget exhausted\\
2 & 4 & 2 & $[0,3]$ & Budget exhausted\\
4 & 7 & 4 & $[0,3]$ & Budget exhausted\\
11 & 4 & 11 & $[3,2]$ & Budget exhausted\\
20 & 4 & 20 & $[5,1]$ & Budget exhausted\\
29 & 4 & 29 & $[6,0]$ & Budget exhausted\\
31 & 16 & 31 & $[6,0]$ & Budget exhausted\\
32 & 17 & 32 & $[6]$ & Budget exhausted\\
33 & 18 & 33 & $[6]$ & Halt\\
\bottomrule
\end{tabular}
\caption{Fresh invocations of the valid countdown image with sufficient capacity. Equal PCs can have different state and retirement counts. Reaching HALT after consuming the final budget unit returns halt, not budget exhaustion. These are guest transitions, not native instruction counts.}
\label{tab:revision-bytecode-prefixes}
\end{table}


\Needspace{8\baselineskip}
The native unit bundles the default countdown. For additional manual experiments,
use the fixed fixture helper from the book checkout after implementing your entry. It edits bytecode
and caller controls while preserving the native instruction text:
\Needspace{6\baselineskip}
\begin{verbatim}
python3 tools/revision_vm_programs.py countdown /tmp/vm-four.asm \
  --source /path/to/your/independent/vm.asm --budget 4
\end{verbatim}
Open the generated copy and break at \texttt{vm\_ready}. The descriptor already
contains the return state; the caller has not stored \texttt{last\_status} and
\texttt{last\_steps} yet. Expect PC seven, depth two, retired four, status zero,
fault minus one and live values zero and three. The helper refuses existing
outputs and writes a receipt naming the selected image and controls. It
supplies no interpreter implementation.

\begin{figure}[!ht]
\centering
\begin{tikzpicture}[alt={Actual native descriptor after countdown with guest budget four. First-row words are PC seven, depth two, retired four and status zero. Native outlines distinguish the PC and retired count. The second row reports fault minus one and preserves inputs.}]
\node[anchor=north west,inner sep=0,font=\small\bfseries] (title) at (0,0) {PC seven is not a count of seven instructions};
\node[anchor=north west,inner sep=0] (a) at ([yshift=-2mm]title.south west) {\MARSRevisionCrop{revision/captures/bytecode/vm-budget.png}{0.99\linewidth}{0.03125}{0.31555556}{0.58203125}{0.60888889}};
\begin{scope}[shift={(a.south west)},x={($(a.south east)-(a.south west)$)},y={($(a.north west)-(a.south west)$)}]
\draw[black,line width=.7pt] (.20,.3088) rectangle (.39,.6176);\draw[black,line width=.7pt] (.58,.3088) rectangle (.77,.6176);
\end{scope}
\end{tikzpicture}
\caption{After four successful guest instructions, PC is seven and depth is two. Solid outlines distinguish PC from retired count four. Status zero and fault minus one identify budget exhaustion; live values zero and three are checked separately. Native LaTeX supplies the title, crop and outlines.}
\label{fig:revision-bytecode-native-budget}
\end{figure}


\subsection*{Prove stack requirements before committing an operation}

For a binary operation, depth must be at least two before either operand is
loaded. The older value is left and top is right. Thus pushing seven, then
two, then SUB produces five, not minus five. Load each as a signed word, form
a wide candidate, and check signed-word range before replacing the pair.
The arithmetic reasoning from the expression lab applies, but this VM uses
its own status codes.

Division first requires both operands, then rejects a zero right operand,
then checks the quotient range. Its quotient truncates toward zero. Minimum
signed word divided by minus one is arithmetic overflow. A failure leaves the
pre-instruction live prefix and depth intact; it also leaves PC and retired
count at their pre-instruction values. Inactive cells within capacity are
unspecified workspace, so preserving the live prefix does not mean clearing
or comparing every physical stack word.

DUP needs one value and room for another. OVER needs two and room for another.
Check underflow before room. With an empty stack and capacity zero, DUP reports
underflow, not overflow. With the countdown, capacity two and budget two produce
an ordinary budget return at PC four with \([0,3]\). Budget three instead attempts
DUP and reports stack overflow: status minus six, PC four, depth two, retired
two, fault four, with that same live prefix. The failed DUP has not retired.

Use a fresh \texttt{countdown} fixture with \texttt{--capacity 2 --budget 3}
to reproduce that case. The native figures show a completed interpreter at the
same public caller boundary; your empty skeleton cannot yet produce them.
Private handler registers and frames are not required to match those figures.

\begin{figure}[!ht]
  \centering
  \begin{tikzpicture}[alt={Actual native descriptor after DUP fails for lack of room. First-row words are PC four, depth two, retired two and stack-overflow status minus six. The second row contains capacity two and fault four. Native outlines identify returned execution state and the capacity/fault pair.}]
    \node[anchor=north west,inner sep=0,font=\small\bfseries] (title) at (0,0) {The failed DUP does not retire};
    \node[anchor=north west,inner sep=0] (a) at ([yshift=-2mm]title.south west) {\MARSRevisionCrop{revision/captures/bytecode/vm-capacity.png}{0.99\linewidth}{0.03125}{0.31555556}{0.58203125}{0.60888889}};
    \begin{scope}[shift={(a.south west)},x={($(a.south east)-(a.south west)$)},y={($(a.north west)-(a.south west)$)}]
      \draw[black,line width=.7pt] (.20,.3088) rectangle (.965,.6176);\draw[black,line width=.7pt,dashed] (.20,.005) rectangle (.58,.303);
    \end{scope}
  \end{tikzpicture}
  \caption{Capacity two and budget three leave PC four, depth two and retired count two. Status minus six reports stack overflow; fault is four. The lower dashed outline includes preserved input capacity two and output fault four. Separate checks confirm the pre-failure live prefix remains zero and three. Native LaTeX supplies the title, crop and outlines.}
  \label{fig:revision-bytecode-native-capacity}
\end{figure}


For \texttt{add-overflow}, the helper generates PUSH32 2147483647, PUSH8 1,
ADD and HALT. ADD fails at PC seven after two retirements, preserving both
operands. For \texttt{zero-divisor}, it generates PUSH8 7, PUSH8 0, DIV and
HALT; DIV fails at PC four after two retirements. Neither failure is a license
to execute the later HALT or report its success.

\subsection*{Retire successful instructions and check budget before fetch}

Maintain a pending next PC until the instruction succeeds. Sequential operations
add their encoded byte size; JMP selects its target; JZ chooses target or PC+2
and pops its condition on both paths. Successful retirement installs that next
PC, increments the retired count and consumes one guest budget unit. HALT also
advances and retires once, then returns halt immediately.

Budget exhaustion is checked before the next fetch. For the two-byte image
PUSH8 7, budget one returns status zero at PC two with one retired instruction.
Budget two instead attempts another fetch and reports fetch-at-end, status minus
four, PC two, retired one. The first instruction's state remains visible.
An empty image with zero budget returns budget exhaustion; with positive budget
it reports fetch at the image end. A valid one-byte HALT with budget one returns
halt, not budget exhaustion.

\Needspace{17\baselineskip}
Do the whole-image validation even when budget is zero. The guest count measures
semantic transitions, not the MIPS work needed to establish that they are safe.
The public measurement helper runs a known countdown image at budgets zero and
64 and counts the complete \texttt{run\_vm} call:
\Needspace{6\baselineskip}
\begin{verbatim}
java -jar /path/to/cli.jar grade \
  independent/work.kts \
  /path/to/your/independent/vm.asm
\end{verbatim}

\begin{table}[!ht]
\centering
\begin{tabular}{r r r r}
\toprule
Guest budget & Retired guest instructions & Native, bitmap design & Native, byte-map design\\
\midrule
0 & 0 & 734 & 599\\
64 & 33 & 2222 & 2087\\
\bottomrule
\end{tabular}
\caption{Measured complete-function work in two qualified implementations. The byte-map implementation also groups opcode classes, so the difference cannot be assigned solely to boundary representation. Setup, assembly, session creation and exit are excluded. These are simulated MIPS instruction counts, not elapsed time or hardware cycles.}
\label{tab:revision-bytecode-work}
\end{table}


These measured designs both satisfy the same behavior. The helper reproduces
the measurement procedure on your implementation; it need not reproduce their
exact counts. Grade correctness first. A faster routine that skips validation
at zero budget is implementing a different contract. A self-branch can retire
the same guest instruction repeatedly until budget exhaustion, even though its
PC never advances to a new byte.

\subsection*{Publish the result and test fresh invocations}

Return status in \texttt{v0} and retired guest count in \texttt{v1}. Every return
also writes PC, depth, retired count, status and fault in the descriptor. This
includes validation failure and zero budget. Capacity and the full stack pointer
are preserved inputs. Unlike a decoder that promises untouched output on
failure, this entry must publish a complete failure report.

Statuses zero and one mean budget exhaustion and halt; fault is minus one in
both. Validation failures are minus one for opcode, minus two for truncation
and minus three for target; they publish zero PC, depth and retired count,
with the offending instruction start as fault. Runtime failures are minus four
for fetch at end, minus five for underflow, minus six for stack overflow,
minus seven for arithmetic overflow and minus eight for zero division. They
publish the failing PC as fault and retain the preceding live stack state.
Inspect status before interpreting a negative value as either a guest operand
or an error code.

All caller access must fit the selected code slice, the 32-byte descriptor or
the selected stack capacity. Code is read-only. Only the five output fields
and selected stack capacity may be written, including temporary writes.
Guards, descriptor inputs, other arenas and cells outside capacity are preserved.
Use full addresses when adding guest offsets or stack indices. Preserve every
bit of the public saved registers, frame and stack pointers, and return through
the incoming native return address. Submission-owned constants and private
native frames are allowed. Produce no console output.

Start construction with record classification and boundary marking, then add
target validation and complete publication of errors. Establish budget/fetch
order before implementing handlers. Make push and simple stack operations work,
then checked arithmetic and the two branches. Keep the common retirement path
separate from failures so adding one handler does not accidentally charge a
failed instruction. These are design stages rather than required private labels
or a fixed frame layout.

\Needspace{7\baselineskip}
Use the six native assessments as construction checkpoints on the same source.
The structural check retains 500 original opcode/truncation failures; the target
check retains eighteen invalid-target cases. The completed-outcome check contains
421 halts or budget returns, and the runtime check contains 308 failures. Together
these are all 1,247 original direct-call trials, with their selected arenas,
offsets, capacities, budgets and saved-state sentinels retained.

Four new public replacement sequences each make three calls in one machine.
The caller check runs twenty-three complete programs, observing initialization
at \texttt{vm\_call}, the actual VM return at \texttt{vm\_ready}, and the later
\texttt{last\_status/last\_steps} stores at \texttt{vm\_stored}. Your entry must
publish the descriptor before return; it must leave those two caller-owned words
for the caller to store. The final interval also protects data after publication.

Debug the default completed caller at those three boundaries. Before the call,
expect length eighteen and guest budget sixty-four. At return, expect
\texttt{v0}=1 and \texttt{v1}=33. Inspect PC eighteen, depth one, retired count
thirty-three, status one and fault minus one in the descriptor, and value six in
the live stack prefix. Continue to \texttt{vm\_stored} and inspect the two stored
words before the exit selector changes \texttt{v0}. The native figures above
remain observations of the supplied public interface in a completed interpreter;
your private handler layout need not match it.

From the imported package directory, the optional command-line route is:
\begin{verbatim}
java -jar /path/to/cli.jar grade \\
  independent/kernel.kts independent/vm.asm
\end{verbatim}
Repeat with \texttt{independent/caller.kts}. The extra public replacement sequences
run in the native lesson. The unfinished skeleton deliberately fails. Equivalent
decoders, boundary sets and dispatch organization are accepted; exact native
instruction counts and inactive workspace contents are not graded. Finish with
ordinary Run, which should exit silently; that transcript alone cannot establish
correct guest execution.

Finally, replace an image between calls without changing its address or length.
An all-HALT image marks every byte as a start. A replacement containing operands
must not inherit those markers. Reinitialize private boundary storage and all
logical execution state from the new invocation. Reusing the old boundary map
can make a branch into an operand appear valid; resuming old output PC can skip
new instructions. Test those possibilities after both a success and an error.
Retain your byte-boundary map, a guest-state trace and an explanation of each
return field alongside your working implementation.
